\documentclass{article}
\input{../../../lib.sty}

\title{\texttt{fn get\_margin}}
\author{Michael Shoemate}
\date{}

\begin{document}

\maketitle

\contrib
Proves soundness of \rustdoc{domains/polars/frame/fn}{get\_margin} in \asOfCommit{mod.rs}{f5bb719}.

\texttt{get\_margin} returns a \texttt{Margin} for a given set of \texttt{grouping\_columns} 
whose descriptors are no more restrictive than what is known in \texttt{FrameDomain}.

\section{Hoare Triple}
\subsection*{Precondition}
\subsubsection*{Compiler-verified}
\begin{itemize}
    \item Argument \texttt{domain} is of type \texttt{FrameDomain}
    \item Argument \texttt{grouping\_column} is of type \texttt{BTreeSet<String>}
\end{itemize}

\subsubsection*{Human-verified}
None

\subsection*{Pseudocode}
\lstinputlisting[language=Python,firstline=2,escapechar=|]{./pseudocode/get_margin.py}

\subsection*{Postcondition}
Returns a \texttt{Margin} that describes properties of members of \texttt{domain} when grouped by \texttt{grouping\_columns}.

\section{Proofs}

\begin{proof} 
    On line \ref{let-margin}, \texttt{margin} is either a valid margin descriptor 
    for \texttt{grouping\_columns} by the definition of \texttt{domain},
    or it is the default margin, which is a valid margin descriptor for all potential datasets. 

    We now update descriptors based on other information available in \texttt{domain}.

    On line \ref{let-subset-margins}, 
    \texttt{subset\_margins} is the subset of margins spanned by \texttt{grouping\_columns}.
    Then \ref{let-mpl} and \ref{let-mpc} assign the smallest known descriptors over any margin spanning a subset of the grouping columns.

    If \texttt{max\_partition\_length} or \texttt{max\_partition\_contributions} 
    is known about a coarser data grouping (when grouped by fewer columns),
    then these descriptors still apply to a finer data grouping, as partition length or per-partition contributions 
    can only decrease when more finely splitting data.

    Therefore \texttt{max\_partition\_length} and \texttt{max\_partition\_contributions} remain valid after mutation.

    Line \ref{all-mnps} retrieves all known \texttt{max\_num\_partitions} descriptors.
    There are no manual preconditions to \rustdoc{domains/polars/frame/fn}{find\_min\_covering}, 
    therefore we claim the postcondition, that the output is a covering for \texttt{grouping\_columns}.

    The number of partitions can be no greater than the cardinality of the cartesian product of the partitions for each of the grouping keys.
    Therefore the code finds a set of \texttt{max\_num\_partitions} descriptors that covers the grouping columns,
    and then updates \texttt{margin} to the product of the covering.

    On an aside, for utility, while the covering found may not be the smallest, 
    the greedy algorithm will always choose a singleton cover if it is available,
    therefore this update to the descriptor cannot increase the descriptor.

    The same logic applies when updating the descriptor for \texttt{max\_influenced\_partitions}.
    Therefore \texttt{max\_num\_partitions} and \texttt{max\_influenced\_partitions} remain valid after mutation.

    Finally, on \ref{all-infos}, \texttt{all\_infos} contains descriptors for grouping key invariants for any margin that includes \texttt{grouping\_columns}.
    If partition keys and/or lengths are known for a finer partitioning, then they are also known for a coarser partitioning.
    Therefore \texttt{public\_info} is updated to the most permissive descriptor for a partitining as fine or finer than \texttt{grouping\_columns}.

    Since the initial margin (\ref{let-margin}) was valid, and all updates have also been shown to be valid,
    \texttt{get\_margin} returns a \texttt{Margin} that describes properties of members of \texttt{domain} when grouped by \texttt{grouping\_columns}.
\end{proof}
\end{document}
